% Brief (tikz-diagram-craft), register row ch4-48.
% Reader question: as each readers critical set grows, how much worse than
%   just doubling the single-reader floor does running two separate digest
%   heads actually get?
% Claim: the split penalty Floor(N,2k,m)/(2 Floor(N,k,m)) is always above 1
%   and grows without bound in k, not a fixed band; at N=60,m=8 it runs
%   1.029 to 1.224 over k=1..4, and at N=60,m=20 it reaches 1.392 by k=10.
% Evidence: penalty(60,k,8) for k=1..4 = 1.029,1.067,1.122,1.224; and
%   penalty(60,10,20)=1.392 [verified, closed form, this pass].
% Must distinguish: the two review-budget curves (m=8, m=20); that neither
%   curve ever returns to the naive-doubling ratio of 1.
% Grammar chosen: measured-quantity plot, penalty ratio against k, two
%   review-budget curves (m=8, m=20), a naive-doubling guide at ratio=1.
% Grammar rejected: a bar chart at the named k values would state the
%   numbers without showing the trend is monotone and unbounded in k, which
%   is the theorems own content, not a property of five named points alone.
% Five-second test: does the m=20 curve ever dip back toward the naive
%   doubling line as k grows, or only move further from it?
\begin{figure}[H]
\centering
\pdfigurehue{pdteal}%
\begin{tikzpicture}[pd figure]
\begin{axis}[pd axis,
  width=7.8cm,height=5.2cm,
  xmin=0.7,xmax=11.3,ymin=0.98,ymax=1.43,
  xtick={1,2,3,4,5,6,7,8,9,10},
  ytick={1.0,1.1,1.2,1.3,1.4},
  xlabel={critical-set size $k$},
  ylabel={split penalty ratio}]
  \addplot[pd warn series,dashed,line width=.9pt,domain=0.7:11.3,samples=2] {1};
  \node[pd label,anchor=south west] at (axis cs:7,1.0) {naive $2\times$: ratio $=1$};
  \addplot[pd focus rule,line width=1.05pt,mark=*,mark size=1.5pt,mark options={fill=pdfocus}]
    coordinates {(1,1.0290)(2,1.0670)(3,1.1221)(4,1.2240)};
  \addplot[pd rule,line width=1.05pt,mark=*,mark size=1.5pt,mark options={fill=pdink}]
    coordinates {(1,1.0157)(2,1.0332)(3,1.0529)(4,1.0754)(5,1.1015)(6,1.1323)
                 (7,1.1698)(8,1.2175)(9,1.2830)(10,1.3922)};
  \node[pd focus label,anchor=south east] at (axis cs:4,1.224) {$m{=}8$};
  \node[pd label,anchor=west] at (axis cs:10.15,1.3922) {$m{=}20$};
  \node[pd focus label,anchor=north west] at (axis cs:4.15,1.20) {1.224};
  \node[pd label,anchor=south] at (axis cs:9.5,1.40) {1.392};
\end{axis}
\end{tikzpicture}
\caption{The super-additive split penalty of Theorem~\ref{thm:split-penalty},
$\mathrm{Floor}(N,2k,m)/(2\,\mathrm{Floor}(N,k,m))$ at $N{=}60$, against
critical-set size $k$, for two review budgets. Both curves stay strictly
above the naive-doubling ratio of $1$ (dashed) and grow without bound in
$k$: at $m{=}8$ the penalty runs from $1.029$ at $k{=}1$ to $1.224$ at
$k{=}4$ (its domain ends there, since $2k\le m$), and at $m{=}20$ it
reaches $1.392$ by $k{=}10$. Running two digest heads for disjoint readers
therefore costs strictly more than twice one head, and increasingly more
as each readers critical set grows. [verified, closed form]}
\label{fig:split-penalty}
\end{figure}
